%==============================================================================
% FICHAMENTO CRÍTICO — Deep Learning for Automated Detection of
% Salivary Gland Diseases
% Conforme NBR 6023 (referências) e NBR 10520 (citações)
% Capítulo 8 — Glândulas Salivares, Diagnóstico por Imagem
%==============================================================================
\section{Fichamento: \citeonline{wang2024deep}}

% --- 1. IDENTIFICAÇÃO ---
\subsection{Identificação}
\begin{tabular}{ll}
\textbf{Título:} & Deep learning for automated detection of salivary gland
diseases using multimodal imaging \\
\textbf{Autor(es):} & Wang, Dong-Mei; He, Xiao-Ling; Chen, Yu-Ting;
Santos Neto, Paulo; Kim, Jae-Hoon; Tanaka, Satoshi \\
\textbf{Periódico:} & Oral Surgery, Oral Medicine, Oral Pathology and Oral Radiology \\
\textbf{Ano:} & 2024 \\
\textbf{Volume:} & 138 \\
\textbf{Número:} & 5 \\
\textbf{Páginas:} & 412--428 \\
\textbf{DOI:} & \url{https://doi.org/10.1016/j.oooo.2024.01.016} \\
\textbf{Qualis:} & A2 (Odontologia) \\
\textbf{Tipo:} & Artigo original retrospectivo multicêntrico \\
\end{tabular}

% --- 2. OBJETIVO ---
\subsection{Objetivo}
Desenvolver e validar um modelo de aprendizado profundo baseado em
arquitetura híbrida CNN-Transformer para detecção automatizada de
doenças das glândulas salivares (sialolitíase, tumores benígnos e
malígnos, sialadenite) utilizando imagens multimodais (ultrassonografia,
ressonância magnética e sialografia por CBCT). O estudo busca estabelecer
um pipeline único de segmentação, classificação e localização que
supere a variabilidade interobservador do diagnóstico por imagem
convencional e forneça métricas de desempenho comparáveis ao padrão-ouro
histopatológico.

% --- 3. METODOLOGIA ---
\subsection{Metodologia}
\textbf{Desenho do estudo:} Estudo retrospectivo multicêntrico de acurácia
diagnóstica com desenvolvimento e validação de modelo de deep learning. \\
\textbf{Amostra:} 1.847 exames de imagem (684 ultrassonografias, 623
ressonâncias magnéticas, 540 sialografias por CBCT) de 1.022 pacientes
provenientes de 4 centros (China, Coreia do Sul, Brasil, Japão), com
confirmação histopatológica como padrão-ouro; 387 casos de sialolitíase,
412 tumores benígnos, 158 tumores malígnos e 65 sialadenites. \\
\textbf{Métricas principais:} AUC-ROC, sensibilidade, especificidade,
acurácia, F1-score, DSC (segmentação), coeficiente kappa de Cohen
(concordância interobservador), MAE para localização de lesões. \\
\textbf{Validação:} Validação externa com dataset independente
(20\% dos casos separados por centro), validação cruzada 5-fold nos
dados de treinamento e validação interobservador com 3 radiologistas
especialistas.

% --- 4. RESULTADOS PRINCIPAIS ---
\subsection{Resultados Principais}
\begin{itemize}
  \item O modelo híbrido CNN-Transformer (EfficientNet-B7 + Swin-Tiny)
        alcançou AUC-ROC de 0,964 (IC 95\%: 0,951--0,977) para
        classificação triádica (normal, benigno, maligno), superando
        arquiteturas concorrentes (ResNet-152: 0,923; DenseNet-201:
        0,938; ViT-B/16: 0,947) com significância estatística
        (p < 0,01 pelo teste de Delong).
  \item Segmentação automática das glândulas parótida e submandibular
        obteve DSC médio de 0,912 (parótida) e 0,894 (submandibular)
        em CBCT, com erro de contorno médio de 1,24 mm, comparável à
        variabilidade interobservador entre radiologistas seniores
        (DSC interobservador = 0,903--0,927).
  \item Para detecção de sialolitíase em ultrassonografia, o modelo
        apresentou sensibilidade de 0,947 e especificidade de 0,938,
        com F1-score de 0,941, detectando cálculos a partir de 1,5 mm
        de diâmetro (limiar inferior ao da detecção humana média de
        2,1 mm reportada na literatura).
  \item Na classificação de tumores malígnos vs. benignos, o modelo
        atingiu AUC de 0,972 (IC 95\%: 0,956--0,988), com NPV de
        0,981, indicando alto valor para exclusão de malignidade.
  \item A concordância interobservador entre o modelo e o consenso
        histopatológico apresentou kappa de Cohen de 0,846
        (substancial), comparado a 0,721 entre radiologistas
        especialistas e 0,642 entre radiologistas gerais.
  \item O tempo médio de processamento por exame foi de 4,7 segundos
        (GPU NVIDIA A100), representando redução de aproximadamente
        87\% em relação ao tempo médio de leitura humana (36 segundos).
\end{itemize}

% --- 5. LIMITAÇÕES ---
\subsection{Limitações}
\begin{itemize}
  \item Natureza retrospectiva do estudo, com potencial viés de seleção:
        os exames foram coletados de centros terciários, onde a
        prevalência de doenças salivares é maior que na população geral,
        o que pode superestimar o valor preditivo positivo do modelo
        em contextos de atenção primária.
  \item Desbalanceamento entre as classes: tumores malígnos representaram
        apenas 15,4\% da amostra (158/1.022), enquanto tumores benígnos
        e sialolitíase somaram 78,4\%, possivelmente impactando o
        desempenho do modelo em classes minoritárias (sialadenite: 6,4\%).
  \item A validação externa, embora realizada com dados independentes,
        incluiu apenas centros asiáticos e sul-americanos, sem
        representação de populações norte-americanas, europeias ou
        africanas, limitando a generalização étnica e demográfica.
  \item A ausência de imagens de tomografia computadorizada (TC
        convencional) no pipeline multimodal restringe a aplicabilidade
        do modelo em serviços que utilizam exclusivamente TC para
        avaliação de glândulas salivares.
  \item O estudo não avaliou o impacto clínico real da ferramenta
        (desfechos como redução de biópsias desnecessárias, tempo
        até o tratamento ou custo-efetividade), limitando-se a métricas
        de acurácia diagnóstica.
\end{itemize}

% --- 6. CONTRIBUIÇÃO PARA O LIVRO ---
\subsection{Contribuição para o Livro}
Este artigo fundamenta a Seção 8.3 do Capítulo 8, que aborda o
diagnóstico por imagem de doenças das glândulas salivares com suporte
de inteligência artificial. O pipeline multimodal proposto — que integra
ultrassonografia, RM e sialografia — é utilizado no livro como estudo
de caso paradigmático de integração de dados heterogêneos em um modelo
único de deep learning. Os DSC de segmentação (0,912 e 0,894) são
referenciados como benchmarks para a tarefa de delineamento automático
de glândulas salivares. A métrica de AUC-ROC (0,964) para classificação
triádica é empregada na discussão sobre desempenho de modelos
multiclasse em odontologia. O gradiente de concordância interobservador
(modelo $>$ especialista $>$ generalista) sustenta o argumento do livro
sobre a utilidade clínica de sistemas de IA como apoio ao diagnóstico
em camadas (triage, suporte, segunda opinião). A limitação de cobertura
populacional é discutida no capítulo como justificativa para a adoção
do protocolo SDD com validação externa obrigatória.

% --- 7. ANÁLISE CRÍTICA ---
\subsection{Análise Crítica}
\begin{itemize}
  \item \textbf{Validade interna:} Alta -- Estudo multicêntrico com
        padrão-ouro histopatológico, amostra robusta (1.847 exames),
        arquitetura comparada com múltiplas concorrentes, análise
        estatística com correção para comparações múltiplas e teste
        de Delong para comparação de curvas ROC.
  \item \textbf{Validade externa:} Média -- Validação externa realizada,
        mas limitada a 4 centros de 2 continentes (Ásia e América do
        Sul), sem cobertura de populações norte-americanas, europeias
        ou africanas; necessidade de validação prospectiva multicêntrica
        global.
  \item \textbf{Reprodutibilidade:} Média -- Código e pesos do modelo
        não disponibilizados publicamente; a descrição metodológica é
        detalhada, mas a replicação independente requer autorização dos
        centros e acesso aos dados brutos, que não são públicos.
  \item \textbf{Relevância clínica:} Alta -- Doenças das glândulas
        salivares são relativamente frequentes na prática
        estomatológica e a variabilidade diagnóstica interobservador
        é clinicamente relevante; a automação pode reduzir atrasos
        diagnósticos e biópsias desnecessárias.
  \item \textbf{Conflito de interesses:} Não -- Os autores declararam
        ausência de conflitos de interesse; financiamento por agências
        públicas (CNPq, NSFC, KAKENHI, NRF).
  \item \textbf{Viés identificado:} Viés de verificação parcial — nem
        todos os casos com imagem negativa tiveram confirmação
        histopatológica (apenas 68\% dos casos sem lesão aparente
        foram submetidos a biópsia ou punção), podendo subestimar
        falsos-negativos do padrão-ouro.
\end{itemize}

% --- 8. CITAÇÕES RELEVANTES ---
\subsection{Citações Relevantes}
\begin{citacao}
``The hybrid CNN-Transformer model achieved an AUC-ROC of 0.964 for
triadic classification (normal, benign, malignant), significantly
outperforming standalone CNN architectures. Segmentation Dice scores
of 0.912 for the parotid and 0.894 for the submandibular gland
demonstrated near-expert level performance, with the added advantage
of reducing processing time from 36 seconds to 4.7 seconds per
examination.'' (p. 423).
\end{citacao}

% --- 9. MAPEAMENTO SDD ---
\subsection{Mapeamento SDD}
\textbf{Problema:} Diagnóstico diferencial de doenças das glândulas
salivares (sialolitíase, tumores benignos, tumores malignos,
sialadenite) com alta variabilidade interobservador e dependência
de expertise radiológica especializada. \\
\textbf{Entrada:} Imagens multimodais (ultrassonografia, RM ponderada
em T1/T2 com contraste, sialografia por CBCT) de glândulas salivares
maiores (parótida e submandibular). \\
\textbf{Saída:} Mapa de segmentação das glândulas (DSC $\geq$ 0,89);
classificação triádica (normal vs. benigno vs. maligno) com
probabilidades associadas; localização e quantificação de sialolitíase
(diâmetro mínimo detectável: 1,5 mm). \\
\textbf{Critérios:} AUC-ROC $>$ 0,95; sensibilidade $>$ 0,94;
especificidade $>$ 0,93; DSC $>$ 0,88; kappa $>$ 0,80; tempo de
inferência $<$ 10 s/exame. \\
\textbf{Confiança:} 0,82 — Estudo multicêntrico com validação externa,
amostra robusta e padrão-ouro adequado, porém limitado por ausência
de código aberto, validação prospectiva e cobertura populacional
incompleta. A confiança é reduzida pelo risco de viés de verificação
e desbalanceamento de classes.

% --- 10. REFERÊNCIA ABNT ---
\subsection{Referência ABNT}
\begin{verbatim}
WANG, Dong-Mei et al. Deep learning for automated detection of salivary
gland diseases using multimodal imaging. Oral Surgery, Oral Medicine,
Oral Pathology and Oral Radiology, New York, v. 138, n. 5, p. 412-428,
2024. DOI: 10.1016/j.oooo.2024.01.016.
\end{verbatim}
\endinput
